% 本文件由 tools/build_book.py 从 dist/reports/assertions.json 生成，请勿手改。
\begin{longtable}{@{}p{0.06\textwidth}p{0.32\textwidth}p{0.22\textwidth}p{0.16\textwidth}@{}}
\caption{章节 3.3 数值断言（生成时间 2026-10-19T16:08:50-0400）}\\
\toprule
类别 & 断言 & 理论值 & 实测值 \\
\midrule
\endfirsthead
\toprule
类别 & 断言 & 理论值 & 实测值 \\
\midrule
\endhead
\bottomrule
\endlastfoot

A & 章节声明对账(4 viz / 依赖 3.1 / 断言 B,C) & docs/\allowbreak{}chapters.yaml 是章节清单的唯一真源，MANIFEST 必须逐项对齐 & 4 个可视化：['energy\_drift\_duel',\allowbreak{} 'phase\_orbit',\allowbreak{} 'modified\_hamiltonian',\allowbreak{} 'long\_term\_orbit']；deps=\allowbreak{}['3.1']；asserts=\allowbreak{}['B',\allowbreak{} 'C']；模块 2 个全部存在 \checkmark \\
A & 全参数角点 $\times$ 时间轴两端的帧体检 & PROTOCOLS 1.1/\allowbreak{}1.3 节：帧里不许出现 NaN/\allowbreak{}Inf；出不了有限值就写占位+掩码+文案 & 64 帧、\allowbreak{}586 个数组、\allowbreak{}117 个 meta 字段全部有限且可打包(\allowbreak{}0 个反例) \checkmark \\
B & 10\textasciicircum{}6 步的招牌帧 & 「跑了 10\textasciicircum{}6 步」必须落在 meta 与数组上，不能靠文案 & n\_steps\_observed=\allowbreak{}1000000，t 末端 50000.0，带宽 6.249999e-04(\allowbreak{}理论 h$^2$/\allowbreak{}4=\allowbreak{}6.250000e-04)，两半之比 1.0000 vs 2.0026；短期误差 RK4 2.61e-06 < Verlet 5.16e-03 \checkmark \\
B & 爆点帧的三件套 & PROTOCOLS 1.3 节：占位值 + 掩码 + 文案；占位值取最后可信值 & h=\allowbreak{}2.5 显式 Euler 在第 70 时刻被掩码截断，掩掉的采样点 86/\allowbreak{}114，占位值 1.694830e+23 \checkmark \\
B & 相空间帧的两个数(闭合 vs 相位) & 辛：点集闭合(\allowbreak{}到不变曲线偏差 =\allowbreak{} 舍入级)；相位照样漂(\allowbreak{}h$^2$/\allowbreak{}24) & Verlet sup=\allowbreak{}1.84e-14(\allowbreak{}舍入级 $\implies$ growth 不读；名义 growth 3.609，det M=\allowbreak{}1.000000000000)、\allowbreak{}相位 -0.250279 圈；RK4 sup=\allowbreak{}2.61e-04(\allowbreak{}growth 2.001，1$-$det=\allowbreak{}1.387e-08)、\allowbreak{}相位 +0.000498 圈 \checkmark \\
B & 修正哈密顿量帧：带宽闭式 + H\_\allowbreak{}h 守恒 & H\_h 的二次型由 2$\times$2 映射给出；det M=\allowbreak{}1 $\iff$ 它精确守恒 & Verlet 带宽 2.500000e-03 =\allowbreak{} 闭式 2.500000e-03(\allowbreak{}比值 1.000000)，H\_h 相对变化 8.86e-14；RK4 的同一公式给出 1.39e-03(\allowbreak{}1$-$det =\allowbreak{} 1.387e-08) \checkmark \\
B & 行星轨道帧：形状 vs 方位 & 形状由 (\allowbreak{}E,\allowbreak{}L) 决定(\allowbreak{}辛性管得了)；方位是相位型误差(\allowbreak{}辛性管不了) & Verlet：能量带 7.414e-04、\allowbreak{}形状 growth 1.000、\allowbreak{}角动量 1.89e-14、\allowbreak{}进动 -6.512e-04 rad/\allowbreak{}圈；RK4：能量带 4.404e-06(\allowbreak{}更小)、\allowbreak{}形状 growth 1.795、\allowbreak{}角动量 4.83e-07、\allowbreak{}进动 +1.761e-07 rad/\allowbreak{}圈 \checkmark \\
B & 时间轴(log10 观测步数)的口径 & t=\allowbreak{}0 $\to$ 1 步、\allowbreak{}t=\allowbreak{}6 $\to$ 10\textasciicircum{}6 步；数组只覆盖已观测的那一段(\allowbreak{}不是把全长曲线裁一下) & t=\allowbreak{}0：1 步 /\allowbreak{} 2 点；t=\allowbreak{}3：1000 步 /\allowbreak{} 1001 点；t=\allowbreak{}6：1000000 步 /\allowbreak{} 1501 点 \checkmark \\
A & 度量函数的闭式对拍 & 线性序列：lo/\allowbreak{}hi/\allowbreak{}amplitude/\allowbreak{}sup/\allowbreak{}slope 与两半斜率全部可手算；纯振荡的趋势/\allowbreak{}带宽 =\allowbreak{} 3/\allowbreak{}(\allowbreak{}$\pi$N) & band=\allowbreak{}[-2.000,\allowbreak{}28.000] amp=\allowbreak{}30.000 slope=\allowbreak{}3.000000000000 两半斜率=\allowbreak{}(\allowbreak{}3.000000000000,\allowbreak{}3.000000000000)；正弦趋势/\allowbreak{}带宽的 3/\allowbreak{}(\allowbreak{}$\pi$N) 律：N=\allowbreak{}2:0.4739(\allowbreak{}理论 0.4775)、\allowbreak{}N=\allowbreak{}4:0.2378(\allowbreak{}理论 0.2387)、\allowbreak{}N=\allowbreak{}8:0.1191(\allowbreak{}理论 0.1194)、\allowbreak{}N=\allowbreak{}16:0.0596(\allowbreak{}理论 0.0597) \checkmark \\
A & power\_\allowbreak{}law 的回收能力 & y =\allowbreak{} C$\cdot$h\textasciicircum{}p $\implies$ 双对数斜率 =\allowbreak{} p、\allowbreak{}截距 =\allowbreak{} C & 拟合 (\allowbreak{}p,\allowbreak{}C)=\allowbreak{}(\allowbreak{}2.000000000000,\allowbreak{}3.000000000000)，对数残差 4.44e-16；舍入地板案例：不剔除 p=\allowbreak{}16.4242、\allowbreak{}剔除后 p=\allowbreak{}2.000000000000 \checkmark \\
A & 相位误差与角度折叠 & $\varphi$ =\allowbreak{} atan2(\allowbreak{}p,\allowbreak{}q)；圈数 =\allowbreak{} [$\varphi$(\allowbreak{}t)$-$$\varphi$(\allowbreak{}0)]/\allowbreak{}2$\pi$(\allowbreak{}含缠绕) & 精确同相 $\to$ 0.000000000000 圈；反向流 $\to$ 6.366197723676 圈(\allowbreak{}理论 2T/\allowbreak{}2$\pi$ =\allowbreak{} 6.366197723676)；wrapped\_angle(\allowbreak{}$\pm$3$\pi$) =\allowbreak{} +$\pi$(\allowbreak{}值域 (\allowbreak{}$-$$\pi$,\allowbreak{}$\pi$]) \checkmark \\
A & 二次不变量的三类映射对拍 & det M =\allowbreak{} 1 $\iff$ Q =\allowbreak{} [[c,\allowbreak{}(\allowbreak{}d$-$a)/\allowbreak{}2],\allowbreak{}[(\allowbreak{}d$-$a)/\allowbreak{}2,\allowbreak{}$-$b]] 满足 M$^{\mathsf T}$QM =\allowbreak{} Q；det$\ne$1 时不存在 & 剪切残差 0.00e+00、\allowbreak{}旋转残差 8.62e-17(\allowbreak{}Q$\propto$$-$I)、\allowbreak{}缩放 symplectic=\allowbreak{}False；辛 Euler 的 Q 沿 2e4 步轨道相对变化 9.16e-15 \checkmark \\
A & 等值线半径的闭式对拍 & r(\allowbreak{}$\theta$) =\allowbreak{} $\surd$(\allowbreak{}2$\cdot$value /\allowbreak{} (\allowbreak{}Q$_0$$_0$cos$^2$$\theta$ + 2Q$_0$$_1$cos$\theta$sin$\theta$ + Q$_1$$_1$sin$^2$$\theta$)) & 椭圆 (\allowbreak{}a,\allowbreak{}b)=\allowbreak{}(\allowbreak{}2.0,\allowbreak{}0.5) 的最大等值线残差 2.22e-16；Verlet h$\to$0 时半径偏差 1.25e-13 \checkmark \\
B & 闭合度对三种轨道的判别 & $\delta$\_k =\allowbreak{} r\_k/\allowbreak{}r\_curve(\allowbreak{}$\theta$\_k) $-$ 1；舍入级 $\implies$ 点集闭合；线性累积 $\implies$ 螺旋(\allowbreak{}前后半段比 $\approx$ 时长比) & 完美圆 sup=\allowbreak{}2.22e-16；螺旋 sup=\allowbreak{}2.857e-01(\allowbreak{}growth 2.134)；舍入抖动 sup=\allowbreak{}4.44e-16(\allowbreak{}at\_roundoff=\allowbreak{}True) \checkmark \\
B & 采样受限的假证据：加密采样就变小 & 轨道是连续曲线、\allowbreak{}采样点离散：「最近回归距离」的下界由采样密度(\allowbreak{}h)决定，与是否闭合无关 & 同一条精确圆：逐圈最近回归距离的均值 h=\allowbreak{}0.1 时 2.509e-02、\allowbreak{}h=\allowbreak{}0.01 时 2.513e-03(\allowbreak{}差 10.0 倍 $\approx$ 步长比，纯粹是采样下界)；而 curve\_closure 两次都是 sup=\allowbreak{}2.22e-16/\allowbreak{}2.22e-16 \checkmark \\
B & 到精确轨道的偏差：有界带宽 vs 线性累积 & 辛方法：偏差 $\le$ 不变带 $\implies$ 有界；非辛：每圈穿出一点 $\implies$ 线性(\allowbreak{}前后半段比 $\approx$ 2) & Verlet sup=\allowbreak{}1.251e-03(\allowbreak{}growth 1.000)；RK4 sup=\allowbreak{}2.614e-04(\allowbreak{}growth 2.001)；末端 RK4/\allowbreak{}Verlet =\allowbreak{} 0.209(\allowbreak{}观测 600 圈) \checkmark \\
B & 度量函数的报错路径 & 量不出定义的量时必须抛错：帧层才有机会写 null/\allowbreak{}占位，而不是让 0 或 NaN 混进结论 & 5 条边界全部抛 ValueError：closure\_stats(\allowbreak{}<4 点)、\allowbreak{}radial\_deviation(\allowbreak{}长度不符)、\allowbreak{}return\_distances(\allowbreak{}周期$\le$0)、\allowbreak{}curve\_closure(\allowbreak{}出发点在 H\_h=\allowbreak{}0 上)、\allowbreak{}power\_law(\allowbreak{}长度不符) \checkmark \\
A & 方法表 + 四个辛方法的 det M = 1 & det M $\equiv$ 1(\allowbreak{}精确到位)且 M$^{\mathsf T}$QM$-$Q 在舍入级 & 最大偏离 1.388e-16；symplectic\_euler: order=\allowbreak{}1 det=\allowbreak{}1.000000000000 resid=\allowbreak{}1.73e-17；verlet: order=\allowbreak{}2 det=\allowbreak{}1.000000000000 resid=\allowbreak{}1.39e-16；leapfrog: order=\allowbreak{}2 det=\allowbreak{}1.000000000000 resid=\allowbreak{}1.39e-16；implicit\_midpoint: order=\allowbreak{}2 det=\allowbreak{}1.000000000000 resid=\allowbreak{}1.74e-17；rk4: order=\allowbreak{}4 det=\allowbreak{}0.999999986128 resid=\allowbreak{}1.39e-08；explicit\_euler: order=\allowbreak{}1 det=\allowbreak{}1.010000000000 resid=\allowbreak{}1.00e-02 \checkmark \\
A & 10\textasciicircum{}6 步的带宽 = 闭式 & 辛 Euler [$-$h/\allowbreak{}(\allowbreak{}2+h),\allowbreak{} h/\allowbreak{}(\allowbreak{}2$-$h)]；Verlet/\allowbreak{}leapfrog [$-$h$^2$/\allowbreak{}4,\allowbreak{} 0](\allowbreak{}h=\allowbreak{}0.05，E$_0$=\allowbreak{}1) & 最大绝对偏差 6.368e-09；symplectic\_euler: 实测 [-2.439023753e-02,\allowbreak{}+2.564102545e-02] 闭式 [-2.439024390e-02,\allowbreak{}+2.564102564e-02] 偏差 6.37e-09；verlet: 实测 [-6.249999456e-04,\allowbreak{}+0.000000000e+00] 闭式 [-6.250000000e-04,\allowbreak{}+0.000000000e+00] 偏差 5.44e-11；leapfrog: 实测 [-6.249999456e-04,\allowbreak{}+0.000000000e+00] 闭式 [-6.250000000e-04,\allowbreak{}+0.000000000e+00] 偏差 5.44e-11 \checkmark \\
A & RK4 / 显式 Euler 的 det M & RK4：det =\allowbreak{} 1 $-$ h$^6$/\allowbreak{}72 + h$^8$/\allowbreak{}576；显式 Euler：det =\allowbreak{} 1 + h$^2$ & RK4 最大相对差 8.004e-09；h=\allowbreak{}0.4: 1$-$det=\allowbreak{}5.575111e-05 vs 5.575111e-05(\allowbreak{}相对差 0.00e+00)；h=\allowbreak{}0.2: 1$-$det=\allowbreak{}8.844444e-07 vs 8.844444e-07(\allowbreak{}相对差 1.26e-10)；h=\allowbreak{}0.1: 1$-$det=\allowbreak{}1.387153e-08 vs 1.387153e-08(\allowbreak{}相对差 8.00e-09)；h=\allowbreak{}0.05: 1$-$det=\allowbreak{}2.169460e-10 vs 2.169460e-10(\allowbreak{}相对差 0.00e+00)；显式 Euler det(\allowbreak{}0.1)=\allowbreak{}1.010000000000 \checkmark \\
A & 单摆精确周期 vs 数值过零点 & T(\allowbreak{}$\theta$0) =\allowbreak{} 4/\allowbreak{}$\surd$(\allowbreak{}1+H) $\cdot$ K(\allowbreak{}k)(\allowbreak{}AGM 算的完全椭圆积分) & 最大相对差 1.667e-09；$\theta$0=\allowbreak{}0.01: AGM 6.283224577 实测 6.283224567(\allowbreak{}相对差 1.67e-09)；$\theta$0=\allowbreak{}0.5: AGM 6.382789698 实测 6.382789687(\allowbreak{}相对差 1.61e-09)；$\theta$0=\allowbreak{}1.0: AGM 6.699975664 实测 6.699975655(\allowbreak{}相对差 1.46e-09)；$\theta$0=\allowbreak{}2.0: AGM 8.349752927 实测 8.349752920(\allowbreak{}相对差 8.09e-10) \checkmark \\
A & 精确流的轨道要素守恒 + 辛方法的形状误差 $\le$ 带宽 & a =\allowbreak{} $-$$\mu$/\allowbreak{}2E、\allowbreak{}e =\allowbreak{} |A|/\allowbreak{}$\mu$、\allowbreak{}A 为 Laplace--Runge--Lenz 矢量；辛方法能量有界 $\implies$ 形状有界 & 精确流最大偏离 2.665e-15；Verlet(\allowbreak{}200 圈，h=\allowbreak{}0.01)形状失真 2.470e-04 对能量带宽 7.414e-04 \checkmark \\
B & Verlet 的能量误差有界(10\textasciicircum{}6 步) & 有界 $\le$ 6.25001e-04 & 0.000625 \checkmark \\
B & RK4 的能量误差线性漂移(10\textasciicircum{}6 步) & 斜率 $\ne$ 0 & $-1.085\times 10^{-7}$ \checkmark \\
B & RK4 的漂移是单向的(单调占比 = 1) & $\le$ 0 & 0 \checkmark \\
B & 有界 vs 漂移的定量判据(两半 sup 之比 + 趋势/带宽) & 辛：后半/\allowbreak{}前半 sup 之比 =\allowbreak{} 1(\allowbreak{}同一带宽)；非辛：=\allowbreak{} 2(\allowbreak{}线性长) & sup 之比：Verlet 1.0000、\allowbreak{}RK4 2.0019；趋势/\allowbreak{}带宽：Verlet 0.005013、\allowbreak{}RK4 1.000000(\allowbreak{}振幅 6.250e-04 /\allowbreak{} 斜率 -1.085e-07，单调占比 1.0000) \checkmark \\
B & Verlet 的角动量守恒(中心力逐位性质) & 漂移 $\le$ 1.00000e-12 & $1.510\times 10^{-14}$ \checkmark \\
B & 角动量 vs 动量：谁是辛性的证据 & 中心力 + 径向 kick $\implies$ L 逐位守恒(\allowbreak{}辛性)；动量守恒只是平移不变性，非辛方法也有 & Kepler 200 圈角动量漂移：Verlet 1.887e-14、\allowbreak{}RK4 4.830e-07(\allowbreak{}差 25589763.1 倍)；LJ 链动量漂移最大 6.017e-14(\allowbreak{}含显式 Euler 6.017e-14) \checkmark \\
B & 尾部采样点的 t = n\_\allowbreak{}steps\_\allowbreak{}done$\cdot$h & $\le$ 0 & 0 \checkmark \\
B & 爆炸停止点的口径一致性 & i\_blowup 是第一个含爆炸状态的采样点，[:i\_blowup] 可信；t 与 idx 都改成 n\_steps\_done & h=\allowbreak{}0.05 显式 Euler 在第 22134 步停止(\allowbreak{}t=\allowbreak{}1106.7)，i\_blowup=\allowbreak{}34/\allowbreak{}1501，爆点 |y|=\allowbreak{}1.000e+12，可信段能量误差 7.184e+23 \checkmark \\
B & 快路径 vs 通用路径、verlet vs leapfrog & 两条路径算的是同一个映射(\allowbreak{}浮点求和顺序也相同)；leapfrog 就是 verlet 的别名 & 21 组对拍最大差 0.0e+00(\allowbreak{}逐位相同)；verlet/\allowbreak{}leapfrog 在 7 个系统上差 0.0 \checkmark \\
B & 修正不变量 H\_\allowbreak{}h 的守恒性 & det M =\allowbreak{} 1 $\iff$ 存在精确二次不变量 Q(\allowbreak{}H\_h 的二次项)；RK4 的 det =\allowbreak{} 1$-$h$^6$/\allowbreak{}72 $\implies$ 不存在 & symplectic\_euler: H\_h 相对变化 9.16e-15(\allowbreak{}真实能量带宽 2.56e-02)；verlet: H\_h 相对变化 1.47e-14(\allowbreak{}真实能量带宽 6.25e-04)；implicit\_midpoint: H\_h 相对变化 4.30e-14(\allowbreak{}真实能量带宽 4.32e-14)；RK4 同一公式给出 4.34e-06(\allowbreak{}=\allowbreak{} 它自己的能量漂移) \checkmark \\
C & RK4 能量漂移率对 h 的幂律 & 5 & 4.99965 \checkmark \\
C & RK4 漂移率的系数 = T/72 & |$\Delta$E/\allowbreak{}E| =\allowbreak{} h$^5$T/\allowbreak{}72(\allowbreak{}每步掉 h$^6$/\allowbreak{}72，共 n=\allowbreak{}T/\allowbreak{}h 步) & 拟合系数 0.693287 vs T/\allowbreak{}72 0.694444(\allowbreak{}相对差 1.666e-03) \checkmark \\
C & Verlet 带宽对 h 的幂律 & 2 & 1.99999 \checkmark \\
C & 辛 Euler 带宽 = h/(2$-$h)(比值检验) & 带宽 =\allowbreak{} h/\allowbreak{}(\allowbreak{}2$-$h)(\allowbreak{}det M=\allowbreak{}1 的推论，见 notes 第 2 节) & 最大相对差 6.479e-05(\allowbreak{}h$\in$[0.1,\allowbreak{}0.00625]，比值 ['0.99993521',\allowbreak{} '0.99999990',\allowbreak{} '0.99999997',\allowbreak{} '0.99999467',\allowbreak{} '0.99997639']) \checkmark \\
C & 谐振子的稳定边界(二分实测) & Verlet/\allowbreak{}辛 Euler：h<2；RK4：h<2$\surd$2$\approx$2.8284 & Verlet 2.000001(\allowbreak{}22 次求值、\allowbreak{}20 次二分)、\allowbreak{}RK4 2.828979(\allowbreak{}14 次求值) \checkmark \\
C & 形状失真的增长比(前/后半段) & 辛方法：有界振荡 $\implies$ growth$\approx$1；RK4：线性漂移 $\implies$ growth$\approx$2 & verlet: sup=\allowbreak{}2.470e-04 growth=\allowbreak{}1.000；symplectic\_euler: sup=\allowbreak{}9.513e-03 growth=\allowbreak{}1.000；implicit\_midpoint: sup=\allowbreak{}2.818e-04 growth=\allowbreak{}1.000；rk4: sup=\allowbreak{}1.833e-07 growth=\allowbreak{}1.795 \checkmark \\
B & 小步长时每步 $\le$5 次迭代 & $\le$ 5 & 4.958 \checkmark \\
C & 隐式中点的迭代次数与 h 的关系 & 迭代解 y =\allowbreak{} y\_n + h/\allowbreak{}2$\cdot$[f(\allowbreak{}y\_n)+f(\allowbreak{}y)] 的收缩率 $\propto$ h $\implies$ 迭代次数随 h 减小 & h=\allowbreak{}[0.5,\allowbreak{} 0.2,\allowbreak{} 0.1,\allowbreak{} 0.05] $\to$ 平均迭代 [11.884,\allowbreak{} 7.021,\allowbreak{} 5.934,\allowbreak{} 4.958]，残差最大 9.992e-15 \checkmark \\
\end{longtable}
